% Generated mechanically from the frozen conventions.tex target.
% Do not edit this generated file; edit the bound locale source instead.

% OK, start here
%

\setcounter{chapter}{1}
\chapter{约定}



\phantomsection
\label{conventions-section-phantom}


\section{说明}
\label{conventions-section-comments}

\noindent
撰写这些文档时所采用约定背后的理念是：选择行得通的约定。

\section{集合论}
\label{conventions-section-sets}

\noindent
我们使用带选择公理的 Zermelo--Fraenkel 集合论。
参见 \cite{Kunen}。我们不使用集合论宇宙（这一点不同于 SGA4）。
我们不会着重讨论集合论问题，但会确保一切都（当然）严格无误，
因而也不会忽略这些问题。


\section{范畴}
\label{conventions-section-categories}

\noindent
一个范畴 $\mathcal{C}$ 由一个对象集以及每一对对象之间的一个态射集组成。
换言之，它就是其他文献中所谓的“小”范畴。
我们也会使用“大”范畴（其对象构成真类的范畴），
但仅限于《范畴》评注 \ref{categories-remark-big-categories} 中列出的那些。

\section{代数}
\label{conventions-section-algebra}

\noindent
在这些笔记中，环均指带有 $1$ 的交换环（即含幺交换环）。
因此，环的范畴以 $\mathbf{Z}$ 为始对象，
以 $\{0\}$ 为终对象（这是唯一满足 $1 = 0$ 的环）。
模均约定为含幺模。参见 \cite{Eisenbud}。

\section{记号}
\label{conventions-section-notation}

\noindent
这里的自然数是 $\mathbf{N} = \{1, 2, 3, \ldots\}$ 的元素。
整数是 $\mathbf{Z} = \{\ldots, -2, -1, 0, 1, 2, \ldots\}$ 的元素。
有理数域记作 $\mathbf{Q}$。
实数域记作 $\mathbf{R}$。
复数域记作 $\mathbf{C}$。
